\section{Limitaciones y trabajo futuro}
\label{sec:limitaciones}

\textbf{Limitaciones conocidas.}
\begin{itemize}
  \item \textbf{Oclusiones.} Si una mano pasa delante de la cara en el modo en vivo, la barba se dibuja encima de la mano.
  \item \textbf{Pelo sobre los ojos.} En el rapado y el \emph{crew cut}, si un flequillo largo tapa un ojo, quedan mechones sueltos, porque no se borra pelo dentro del polígono de la cara para no inventar el ojo.
  \item \textbf{Orejas tapadas.} Si el pelo tapaba las orejas, el rapado pone pelo corto donde irían las orejas, porque no hay información para reconstruirlas.
  \item \textbf{Barbas muy densas.} En barbas reales muy negras y tupidas, el tinte a veces toma solo los bordes: la textura no se distingue dentro de una zona uniforme.
  \item \textbf{Una sola vista en el espejo 3D.} El espejo se construye desde una foto: no tiene nuca, y con giros de más de unos 40° se estira el relieve detrás de la cara.
  \item \textbf{Velocidad en vivo.} El modo en vivo corre a unos 10 cuadros por segundo; con giros rápidos la barba se atrasa un poco.
  \item \textbf{Clasificación de la cara.} Es una heurística calibrada estadísticamente, no un diagnóstico; caras en el límite entre dos formas pueden caer en cualquiera de ellas.
  \item \textbf{Un solo usuario en vivo.} El estado del modo en vivo es único por servidor: dos cámaras a la vez compartirían el anclaje.
\end{itemize}

\textbf{Trabajo futuro.}
\begin{itemize}
  \item Tinte en vivo en el navegador con un \emph{shader} de WebGL y segmentación en JavaScript, para llegar a 30 cuadros por segundo sin servidor.
  \item Optimizar las gaussianas 3D con varias vistas (3DGS entrenado) para completar la cabeza.
  \item Segmentar la barba con un modelo propio, entrenado con las máscaras de CelebAMask-HQ.
  \item Estado del modo en vivo por sesión, para varios usuarios a la vez.
\end{itemize}
